\chapter{Supervised fine-tuning}
\label{cap:sft}

\section{L'obiettivo, in concreto}

L'addestramento supervisionato minimizza l'entropia incrociata~\eqref{eq:ce}
sulle sole posizioni della \emph{completion}, cioè del gloss. Le posizioni del
prompt sono escluse assegnando loro l'etichetta convenzionale $-100$, che le
librerie impiegate interpretano come ``non contribuire alla loss''.

Questa distinzione è più importante di quanto sembri. Se il prompt non venisse
mascherato, il modello dedicherebbe capacità a riprodurre l'inglese di ingresso,
che è già noto e non serve a nulla. Peggio: su questo corpus, dove il gloss è in
larga parte una copia maiuscolizzata del testo, addestrare il modello a
riprodurre il testo rinforzerebbe esattamente la scorciatoia che si vuole
misurare separatamente.

Formalmente, detto $m_t \in \{0,1\}$ l'indicatore ``la posizione $t$ appartiene
alla completion'', la loss effettivamente ottimizzata è
\begin{equation}
  \label{eq:sft-loss}
  \mathcal{L}_{\text{SFT}}(\theta)
  \;=\;
  -\frac{\sum_{i}\sum_{t} m^{(i)}_t
         \log \policy\!\left(y^{(i)}_t \mid x^{(i)}, y^{(i)}_{<t}\right)}
        {\sum_{i}\sum_{t} m^{(i)}_t}.
\end{equation}

\section{Configurazione}

L'implementazione è in \texttt{src/training/sft\_train.py}, funzione
\texttt{run\_sft}, che avvolge \texttt{trl.SFTTrainer}.

\begin{center}
\begin{tabular}{lr}
\toprule
Iperparametro & Valore \\
\midrule
Epoche & $3$ \\
Batch per dispositivo & $4$ \\
Accumulo del gradiente & $4$ \\
Batch effettivo & $16$ \\
Learning rate & $2 \cdot 10^{-5}$ \\
Passi di warmup & $50$ \\
Lunghezza massima di sequenza & $768$ \\
Holdout interno per la validazione & $2\%$ \\
Pazienza per l'arresto anticipato & $3$ \\
\bottomrule
\end{tabular}
\end{center}

Il \emph{warmup} porta il learning rate da zero al valore nominale nei primi
$50$ passi, evitando che i primi aggiornamenti, calcolati su gradienti ancora
disallineati, danneggino i pesi. L'\emph{holdout} del $2\%$ è ritagliato dal
training set (non dal test) e serve solo a decidere l'arresto anticipato: se la
loss di validazione non migliora per $3$ valutazioni consecutive,
l'addestramento si ferma.

Tre epoche su $72\,979$ esempi con batch effettivo $16$ corrispondono a circa
$13\,600$ passi di ottimizzazione. Questo numero va tenuto a mente per il
confronto con il GRPO del capitolo successivo, che ne compie $2\,000$.

\section{Risultati}

Valutazione su $2\,000$ prompt del test set, con campionamento a temperatura
$0{,}7$ e $N=5$ generazioni per prompt (metriche versione 2, si veda il
Capitolo~\ref{cap:metriche}).

\begin{center}
\begin{tabular}{lr}
\toprule
Metrica & SFT \\
\midrule
ROUGE-L & $0{,}9749$ \\
Exact match & $0{,}8117$ \\
Gloss F1 & $0{,}9760$ \\
BLEU (a corpus) & $0{,}9426$ \\
chrF (a corpus) & $97{,}27$ \\
Validity & $0{,}9998$ \\
Non-copy-token accuracy & $0{,}9660$ \\
\bottomrule
\end{tabular}
\end{center}

Nudi, questi numeri sembrano un successo pieno. Vanno letti accanto alla
baseline a regole.

\section{Il confronto che conta: SFT contro regola}

\begin{center}
\begin{tabular}{lrrr}
\toprule
Metrica & Regola & SFT & Delta \\
\midrule
ROUGE-L & $0{,}9685$ & $0{,}9749$ & $+0{,}0064$ \\
Exact match & $0{,}5865$ & $0{,}8117$ & $+0{,}2252$ \\
Non-copy-token accuracy & $0{,}9424$ & $0{,}9660$ & $+0{,}0236$ \\
\bottomrule
\end{tabular}
\end{center}

Gli intervalli di confidenza al $95\%$, ottenuti con bootstrap clusterizzato per
prompt (si veda il Capitolo~\ref{cap:metriche}), sono:

\begin{center}
\begin{tabular}{lc}
\toprule
Delta & Intervallo di confidenza $95\%$ \\
\midrule
ROUGE-L: $+0{,}0064$ & $[+0{,}0038,\; +0{,}0098]$ \\
Exact match: $+0{,}2252$ & $[+0{,}2025,\; +0{,}2510]$ \\
\bottomrule
\end{tabular}
\end{center}

Entrambi gli intervalli escludono lo zero. Va però anticipato subito un
\textbf{caveat che li ridimensiona}, sviluppato nella
sezione~\ref{sec:dispersione-run}: il bootstrap ricampiona i prompt di
valutazione e quindi stima soltanto l'incertezza dovuta al campione di test, non
quella dovuta all'addestramento. La cella SFT, eseguita due volte, varia da sola
di $0{,}0036$ su ROUGE-L e di $0{,}0304$ su exact match. Ne segue che
l'intervallo su ROUGE-L è ottimistico e il delta corrispondente non è
distinguibile dal rumore di addestramento.

Con questa avvertenza, la \emph{grandezza} del vantaggio dipende radicalmente
dalla metrica:

\begin{itemize}
  \item su ROUGE-L il guadagno è $+0{,}0064$, cioè lo $0{,}6\%$ assoluto, dello
        stesso ordine della dispersione fra due run della stessa cella
        ($0{,}0036$). Un modello addestrato su $73$ mila esempi guadagna, su
        questa metrica, quanto il rumore del proprio addestramento;
  \item su exact match il guadagno è $+0{,}2252$, cioè un abbattimento del tasso
        di errore dal $41{,}4\%$ al $18{,}8\%$. In termini di errore relativo è
        una riduzione di circa il $54\%$, ed è circa sette volte la dispersione
        osservata: questo delta regge.
\end{itemize}

Questa discrepanza è la prima manifestazione concreta della saturazione:
ROUGE-L, essendo una metrica di sovrapposizione, non ha risoluzione residua per
distinguere i due sistemi, mentre l'exact match sì. È un argomento operativo per
promuovere l'exact match e la non-copy-token accuracy a metriche primarie
(Capitolo~\ref{cap:metriche}).

\section{Che cosa impara l'SFT che la regola non sa}
\label{sec:sft-vs-regola-analisi}

Il confronto aggregato dice \emph{quanto}, non \emph{che cosa}. È stata quindi
fatta un'analisi caso per caso sui $2\,000$ prompt.

\begin{center}
\begin{tabular}{lr}
\toprule
Esito & Casi \\
\midrule
SFT corretto, regola sbagliata & $586$ \quad ($29{,}3\%$) \\
\quad di cui: cancellazione contestuale della copula \texttt{BE} & $382$ \quad ($65\%$ dei 586) \\
\bottomrule
\end{tabular}
\end{center}

Cioè: due terzi del vantaggio dell'SFT sull'exact match si riducono a una sola
regolarità linguistica. La copula inglese (\emph{is}, \emph{are}, \emph{be})
talvolta compare nel gloss come \texttt{BE} e talvolta viene omessa, e
l'omissione dipende dal contesto. La baseline, essendo context-free, deve
scegliere una volta per tutte e sbaglia sistematicamente sull'altro caso.

\subsection{Il controllo: estendere la regola non funziona}

L'ipotesi naturale è che si possa recuperare quel vantaggio estendendo la
baseline al contesto locale, cioè apprendendo una tabella su bigrammi invece che
su parole singole. È stato provato:

\begin{center}
\begin{tabular}{lr}
\toprule
Baseline & Exact match \\
\midrule
Lessico su unigrammi (originale) & $0{,}5865$ \\
Lessico esteso ai bigrammi & $0{,}5235$ \\
\bottomrule
\end{tabular}
\end{center}

Il risultato \emph{peggiora} di $0{,}063$. La spiegazione è la sparsità: passando
ai bigrammi, molte coppie osservate nel test non compaiono nel training, e la
tabella deve ricadere su un valore di riserva più spesso di quanto guadagni in
precisione contestuale.

Questo controllo è metodologicamente importante perché delimita il contributo
dell'SFT in modo non banale: l'SFT non sta solo memorizzando una tabella più
grande, sta generalizzando una dipendenza contestuale che un conteggio diretto
sui bigrammi non riesce a stimare con i dati disponibili. È il guadagno reale
dell'apprendimento su questo corpus, ed è piccolo ma identificabile.

\section{L'SFT è al tetto del corpus}

Il ROUGE-L dell'SFT è $0{,}9749$. Il miglior valore pubblicato in letteratura su
ASLG-PC12 è $95{,}87$~\cite{peng2023monolingual}, cioè $0{,}9587$.

Segue che l'SFT di un modello da $0{,}5$ miliardi di parametri con adattatori
LoRA \emph{supera} lo stato dell'arte pubblicato su questo corpus. Questo non è
un risultato di cui vantarsi: è la prova più chiara che la metrica non misura
ciò che si crede misuri. Un modello piccolo, addestrato per tre epoche su una
singola GPU, non batte anni di ricerca; batte una metrica satura.

Ne segue anche la conseguenza operativa più importante per il resto della
relazione: \textbf{non c'è spazio di miglioramento residuo che il reinforcement
learning possa colmare}. Con exact match $0{,}8117$ e validity $0{,}9998$, un
metodo di rinforzo partirebbe da una policy già quasi ottima rispetto alla
reward disponibile, in una regione dove i gruppi di candidati tendono a essere
omogenei e quindi il vantaggio~\eqref{eq:advantage} tende a annullarsi. È
esattamente ciò che si osserva, ed è quantificato nel capitolo successivo.
